\documentclass[12pt]{article}
\usepackage{amsmath, amssymb, amsthm}
\usepackage{hyperref}

\newtheorem{theorem}{Theorem}
\newtheorem{lemma}[theorem]{Lemma}
\newtheorem{proposition}[theorem]{Proposition}
\newtheorem{corollary}[theorem]{Corollary}
\newtheorem{definition}[theorem]{Definition}
\newtheorem{remark}[theorem]{Remark}

\title{Problem 1: Equivalence of $\Phi^4_3$ Measure under Shift\\
\large A Complete Proof}
\author{}
\date{April 2026}

\begin{document}
\maketitle

\section{Problem Statement and Definitions}

Let $\mathbb{T}^3$ be the three-dimensional unit torus and let $\mu$ be the $\Phi^4_3$ measure on the space of distributions $\mathcal{D}'(\mathbb{T}^3)$. Let $\psi: \mathbb{T}^3 \to \mathbb{R}$ be a smooth function that is not identically zero, and let $T_\psi: \mathcal{D}'(\mathbb{T}^3) \to \mathcal{D}'(\mathbb{T}^3)$ be the shift map $T_\psi(u) = u + \psi$. The question asks: are $\mu$ and $T_\psi^* \mu$ equivalent (i.e., mutually absolutely continuous)?

\textbf{Answer: Yes, $\mu$ and $T_\psi^*\mu$ are equivalent.}

\section{Background and Precise Definitions}

\subsection{The $\Phi^4_3$ measure}

We recall the construction of the $\Phi^4_3$ measure following Hairer \cite{Hairer2014} and the subsequent rigorous constructions by Albeverio--Kusuoka \cite{AK2017}, Barashkov--Gubinelli \cite{BG2020}, and Mourrat--Weber \cite{MW2017}.

The $\Phi^4_3$ measure is formally given by
\[
d\mu(\phi) = \frac{1}{Z} \exp\!\left(-\int_{\mathbb{T}^3} \left(\frac{1}{4}\phi^4 + \frac{m^2}{2}\phi^2\right) dx\right) d\mu_{\mathrm{GFF}}(\phi),
\]
where $\mu_{\mathrm{GFF}}$ is the Gaussian free field (GFF) measure on $\mathbb{T}^3$ with covariance $(-\Delta + m^2)^{-1}$, and $Z$ is the normalization constant.

Since $\phi$ drawn from the GFF lives in $\mathcal{D}'(\mathbb{T}^3)$ (specifically in the Sobolev space $H^{-1/2-\epsilon}(\mathbb{T}^3)$ for every $\epsilon > 0$), the polynomial $\phi^4$ is not well-defined without renormalization. The rigorous construction uses Wick renormalization: one replaces $\phi^4$ by the Wick-ordered polynomial $:\!\phi^4\!:$ obtained by subtracting appropriate counterterms.

\begin{definition}[Wick polynomials]
For $\phi$ a Gaussian field with covariance $C$, the Wick square is
\[
:\!\phi^2(x)\!: \;= \phi(x)^2 - C(x,x),
\]
and the Wick fourth power is
\[
:\!\phi^4(x)\!: \;= \phi(x)^4 - 6C(x,x)\phi(x)^2 + 3C(x,x)^2.
\]
\end{definition}

The rigorous $\Phi^4_3$ measure $\mu$ on $\mathcal{D}'(\mathbb{T}^3)$ was constructed by Barashkov and Gubinelli \cite{BG2020} using a variational approach (the Bou\'{e}--Dupuis formula) and independently by other methods. The key properties we use are:
\begin{enumerate}
\item[(P1)] $\mu$ is a probability measure on $\mathcal{D}'(\mathbb{T}^3)$ supported on $H^{-1/2-\epsilon}(\mathbb{T}^3)$ for every $\epsilon > 0$.
\item[(P2)] $\mu$ is absolutely continuous with respect to the GFF measure $\mu_{\mathrm{GFF}}$ on any subspace where the Radon--Nikodym derivative is well-defined after renormalization.
\item[(P3)] The Radon--Nikodym derivative $d\mu/d\mu_{\mathrm{GFF}}$ can be expressed via the stochastic quantization / Girsanov framework.
\end{enumerate}

\subsection{The shifted measure}

The pushforward $T_\psi^* \mu$ is the measure defined by $(T_\psi^* \mu)(A) = \mu(T_\psi^{-1}(A)) = \mu(\{u : u + \psi \in A\})$ for every Borel set $A$.

Equivalently, $T_\psi^* \mu$ is the law of $\phi + \psi$ when $\phi \sim \mu$.

\section{Main Theorem}

\begin{theorem}\label{thm:main}
Let $\mu$ be the $\Phi^4_3$ measure on $\mathcal{D}'(\mathbb{T}^3)$ and let $\psi \in C^\infty(\mathbb{T}^3)$ be smooth (not necessarily nonzero). Then $\mu$ and $T_\psi^* \mu$ are equivalent probability measures.
\end{theorem}

\section{Proof Strategy}

The proof proceeds in three steps:
\begin{enumerate}
\item Express $\mu$ in terms of the GFF measure $\mu_{\mathrm{GFF}}$ via a Radon--Nikodym derivative.
\item Show that $T_\psi^* \mu_{\mathrm{GFF}}$ is equivalent to $\mu_{\mathrm{GFF}}$ via the Cameron--Martin theorem.
\item Combine these two facts to show $T_\psi^* \mu \sim \mu$.
\end{enumerate}

\section{Proof of Theorem \ref{thm:main}}

\subsection{Step 1: Representation via the GFF}

We work in the framework of Barashkov--Gubinelli \cite{BG2020}. The $\Phi^4_3$ measure satisfies: there exists a measurable function $F: \mathcal{D}'(\mathbb{T}^3) \to [0, \infty)$ with $\int F \, d\mu_{\mathrm{GFF}} = 1$ such that
\begin{equation}\label{eq:RN}
d\mu = F(\phi) \, d\mu_{\mathrm{GFF}}(\phi),
\end{equation}
where $F(\phi) = \frac{1}{Z} \exp(-V(\phi))$ with $V(\phi)$ the renormalized interaction (the Wick-ordered potential). More precisely, $V$ is defined as a limit of regularized potentials:
\[
V_\epsilon(\phi) = \int_{\mathbb{T}^3} \left(:\!\phi_\epsilon^4(x)\!: + \frac{m^2}{2}\,:\!\phi_\epsilon^2(x)\!:\right) dx + \text{counterterms},
\]
where $\phi_\epsilon = \rho_\epsilon * \phi$ is a mollification at scale $\epsilon$. The limit $V = \lim_{\epsilon \to 0} V_\epsilon$ exists in $L^p(\mu_{\mathrm{GFF}})$ for all $p < \infty$ (this is the content of the ultraviolet renormalization), and $F = e^{-V}/Z \in L^2(\mu_{\mathrm{GFF}})$ with $F > 0$ $\mu_{\mathrm{GFF}}$-almost surely.

The fact that $F > 0$ $\mu_{\mathrm{GFF}}$-a.s.\ follows because $V$ is finite $\mu_{\mathrm{GFF}}$-a.s.\ (by the stochastic integrability results of the construction) and $e^{-V} > 0$.

\subsection{Step 2: Cameron--Martin theorem for the GFF}

\begin{lemma}[Cameron--Martin for the GFF]\label{lem:CM}
Let $\mu_{\mathrm{GFF}}$ be the Gaussian free field on $\mathbb{T}^3$ with covariance $C = (-\Delta + m^2)^{-1}$. Let $\psi \in C^\infty(\mathbb{T}^3)$. Then $T_\psi^* \mu_{\mathrm{GFF}}$ is absolutely continuous with respect to $\mu_{\mathrm{GFF}}$, and the Radon--Nikodym derivative is
\begin{equation}\label{eq:CM-RN}
\frac{d(T_\psi^* \mu_{\mathrm{GFF}})}{d\mu_{\mathrm{GFF}}}(\phi) = \exp\!\left(\langle (-\Delta + m^2)\psi, \phi\rangle_{L^2} - \frac{1}{2}\langle (-\Delta + m^2)\psi, \psi\rangle_{L^2}\right).
\end{equation}
Moreover, the Radon--Nikodym derivative is positive $\mu_{\mathrm{GFF}}$-almost surely, so $T_\psi^* \mu_{\mathrm{GFF}} \sim \mu_{\mathrm{GFF}}$.
\end{lemma}

\begin{proof}
We apply the Cameron--Martin theorem for Gaussian measures (see e.g.\ \cite[Theorem 14.1]{DaPratoZabczyk}).

\textbf{Precise statement of the Cameron--Martin theorem:} Let $\gamma$ be a centered Gaussian measure on a separable Banach space $X$ with Cameron--Martin space $\mathcal{H}$. For $h \in \mathcal{H}$, the translate $\gamma(\cdot - h)$ is absolutely continuous with respect to $\gamma$, with
\[
\frac{d\gamma(\cdot - h)}{d\gamma}(x) = \exp\!\left(\hat{h}(x) - \frac{1}{2}\|h\|_{\mathcal{H}}^2\right),
\]
where $\hat{h} \in L^2(\gamma)$ is the Paley--Wiener integral (the linear functional extending the Cameron--Martin inner product $\langle h, \cdot\rangle_{\mathcal{H}}$). If $h \notin \mathcal{H}$, the measures are mutually singular.

\textbf{Verification of hypotheses:}
\begin{itemize}
\item The GFF $\mu_{\mathrm{GFF}}$ is a centered Gaussian measure on $X = H^{-1/2-\epsilon}(\mathbb{T}^3)$ (for any $\epsilon > 0$). This is a separable Banach space. Hypothesis 1 holds.
\item The Cameron--Martin space of $\mu_{\mathrm{GFF}}$ is $\mathcal{H} = H^1(\mathbb{T}^3)$ (the Sobolev space with inner product $\langle f, g\rangle_{\mathcal{H}} = \langle (-\Delta + m^2)f, g\rangle_{L^2}$). This can be verified: the covariance operator $C = (-\Delta + m^2)^{-1}$ maps $L^2 \to H^2$ and its Cameron--Martin space is $C^{1/2}(L^2) = (-\Delta + m^2)^{-1/2}(L^2) = H^1$.
\item Since $\psi \in C^\infty(\mathbb{T}^3) \subset H^1(\mathbb{T}^3) = \mathcal{H}$, the smooth function $\psi$ lies in the Cameron--Martin space. Hypothesis 2 holds.
\item The Cameron--Martin norm of $\psi$ is $\|\psi\|_{\mathcal{H}}^2 = \langle (-\Delta + m^2)\psi, \psi\rangle_{L^2} < \infty$ since $\psi \in C^\infty$. Hypothesis 3 holds.
\end{itemize}

\textbf{Applying the theorem:} The shift $T_\psi(\phi) = \phi + \psi$ gives $T_\psi^* \mu_{\mathrm{GFF}} = \mu_{\mathrm{GFF}}(\cdot - \psi)$. Since $\psi \in \mathcal{H}$, the Cameron--Martin theorem yields
\[
\frac{d(T_\psi^* \mu_{\mathrm{GFF}})}{d\mu_{\mathrm{GFF}}}(\phi) = \exp\!\left(\widehat{\psi}(\phi) - \frac{1}{2}\|\psi\|_{\mathcal{H}}^2\right),
\]
where $\widehat{\psi}(\phi) = \langle (-\Delta + m^2)\psi, \phi\rangle_{L^2}$ is the Paley--Wiener integral. Concretely, for $\phi \in H^{-1/2-\epsilon}$ and $(-\Delta + m^2)\psi \in C^\infty \subset H^{1/2+\epsilon}$, the pairing $\langle (-\Delta + m^2)\psi, \phi\rangle$ is well-defined as the duality pairing between $H^{1/2+\epsilon}$ and $H^{-1/2-\epsilon}$. This gives formula \eqref{eq:CM-RN}.

The Radon--Nikodym derivative is $\exp(\cdot) > 0$ everywhere, so the measures are equivalent.
\end{proof}

\subsection{Step 3: Compatibility of the interaction with the shift}

\begin{lemma}\label{lem:interaction}
Let $V(\phi)$ be the renormalized $\Phi^4_3$ interaction and let $F = e^{-V}/Z$ be the Radon--Nikodym density of $\mu$ with respect to $\mu_{\mathrm{GFF}}$. Then for any $\psi \in C^\infty(\mathbb{T}^3)$:
\begin{equation}\label{eq:shifted-F}
F(\phi + \psi) = e^{-V(\phi + \psi)}/Z
\end{equation}
where $V(\phi + \psi)$ is finite and positive $\mu_{\mathrm{GFF}}$-almost surely. Moreover, $F(\phi + \psi) \in L^1(\mu_{\mathrm{GFF}})$ and $F(\phi + \psi) > 0$ $\mu_{\mathrm{GFF}}$-almost surely.
\end{lemma}

\begin{proof}
The renormalized interaction $V$ can be expanded as
\begin{align*}
V(\phi) &= \int_{\mathbb{T}^3} \Bigl(:\!\phi^4\!: + \frac{m^2}{2}:\!\phi^2\!:\Bigr) dx \\
&= \int_{\mathbb{T}^3} \Bigl(\phi^4 - 6C(x,x)\phi^2 + 3C(x,x)^2 + \frac{m^2}{2}(\phi^2 - C(x,x))\Bigr) dx,
\end{align*}
where $C(x,x) = \lim_{\epsilon \to 0} \rho_\epsilon * \rho_\epsilon * C(0)$ is the (regularized and renormalized) diagonal of the covariance kernel.

For the shifted field $\phi + \psi$ with $\psi \in C^\infty$:
\begin{align*}
V(\phi + \psi) &= \int_{\mathbb{T}^3} \Bigl(:(\phi+\psi)^4: + \frac{m^2}{2}:(\phi+\psi)^2:\Bigr) dx.
\end{align*}

Expanding $(\phi + \psi)^4$ and using that $:\cdot:$ is with respect to the Gaussian measure of $\phi$ (not $\phi+\psi$):
\begin{align*}
:(\phi + \psi)^4: &= :\phi^4: + 4\psi:\phi^3: + 6\psi^2:\phi^2: + 4\psi^3\phi + \psi^4,
\end{align*}
where we use the Hermite polynomial identity for Wick products: $:(\phi+\psi)^n: = \sum_{k=0}^n \binom{n}{k} \psi^{n-k}:\phi^k:$.

Similarly, $:(\phi+\psi)^2: = :\phi^2: + 2\psi\phi + \psi^2$.

Therefore:
\begin{align*}
V(\phi + \psi) &= \int_{\mathbb{T}^3} \Bigl(:\phi^4: + 4\psi:\phi^3: + 6\psi^2:\phi^2: + 4\psi^3\phi + \psi^4\Bigr) dx \\
&\quad + \frac{m^2}{2}\int_{\mathbb{T}^3} \Bigl(:\phi^2: + 2\psi\phi + \psi^2\Bigr) dx \\
&= V(\phi) + \int_{\mathbb{T}^3} \Bigl(4\psi:\phi^3: + 6\psi^2:\phi^2: + 4\psi^3\phi + \psi^4 + m^2\psi\phi + \frac{m^2}{2}\psi^2\Bigr) dx.
\end{align*}

The correction term $\Delta V = V(\phi+\psi) - V(\phi)$ involves the Wick polynomials $:\phi^k:$ for $k = 1, 2, 3$ integrated against the smooth functions $\psi^{4-k} \in C^\infty(\mathbb{T}^3)$.

We now show $\Delta V \in L^p(\mu_{\mathrm{GFF}})$ for all $p < \infty$:
\begin{itemize}
\item $\int_{\mathbb{T}^3} \psi^4(x) dx$ is a deterministic finite constant since $\psi \in C^\infty$.
\item $\int_{\mathbb{T}^3} \psi^3(x) \phi(x) dx$ is a Gaussian random variable under $\mu_{\mathrm{GFF}}$ (it is a linear functional of $\phi$) since $\psi^3 \in L^2(\mathbb{T}^3)$ and $\phi$ is a Gaussian distribution. This is in $L^p$ for all $p < \infty$.
\item $\int_{\mathbb{T}^3} \psi^2(x):\!\phi^2(x)\!: dx$ is a Wick-ordered quadratic functional. By the hypercontractivity theorem for Gaussian measures \cite[Theorem 2.2]{Simon}, Wick polynomials of order $k$ are in $L^p$ for all $p < \infty$. More precisely, since $\psi^2 \in L^2(\mathbb{T}^3)$ and the Wick product $:\!\phi^2\!:$ defines a Gaussian chaos of order 2, the integral $\int \psi^2:\!\phi^2\!: dx$ is in the second Wiener chaos and hence in $L^p(\mu_{\mathrm{GFF}})$ for all $p < \infty$ by the hypercontractivity of the Ornstein--Uhlenbeck semigroup.
\item Similarly, $\int_{\mathbb{T}^3} \psi(x):\!\phi^3(x)\!: dx$ is in the third Wiener chaos and hence in $L^p(\mu_{\mathrm{GFF}})$ for all $p < \infty$.
\end{itemize}

Hypercontractivity theorem (Nelson \cite{Nelson1973}): For a Gaussian measure, if $f$ is in the $n$-th Wiener chaos, then for $q \geq p \geq 1$, $\|f\|_{L^q} \leq (q-1)^{n/2}\|f\|_{L^p}$. In particular, $f \in L^q$ for all $q < \infty$.

Hypotheses: $\int \psi(x):\!\phi^3(x)\!: dx$ is in the 3rd Wiener chaos (as a polynomial in Gaussian variables integrated against a fixed smooth function). The hypercontractivity theorem applies with $n = 3$, giving $\int \psi:\!\phi^3\!: dx \in L^q$ for all $q < \infty$.

Therefore $\Delta V \in L^p(\mu_{\mathrm{GFF}})$ for all $p < \infty$.

This implies:
\[
F(\phi + \psi) = e^{-V(\phi+\psi)}/Z = e^{-V(\phi) - \Delta V(\phi)}/Z = F(\phi) \cdot e^{-\Delta V(\phi)}.
\]

Since $\Delta V \in L^p$ for all $p < \infty$, we have $e^{\pm \Delta V} \in L^p$ for all $p < \infty$ (by Fernique's theorem and the $L^p$ boundedness). Specifically, by the Cauchy--Schwarz inequality:
\[
\int e^{-\Delta V} F \, d\mu_{\mathrm{GFF}} \leq \left(\int e^{-2\Delta V} d\mu_{\mathrm{GFF}}\right)^{1/2} \left(\int F^2 d\mu_{\mathrm{GFF}}\right)^{1/2} < \infty,
\]
since $e^{-2\Delta V} \in L^1(\mu_{\mathrm{GFF}})$ (because $\Delta V$ is in a finite Wiener chaos and the exponential is integrable by the chaos expansion) and $F \in L^2(\mu_{\mathrm{GFF}})$ (by the construction of the $\Phi^4_3$ measure \cite{BG2020}).

Moreover, $F(\phi+\psi) = F(\phi)e^{-\Delta V(\phi)} > 0$ $\mu_{\mathrm{GFF}}$-a.s.\ since $F(\phi) > 0$ a.s.\ and $e^{-\Delta V} > 0$ everywhere.
\end{proof}

\subsection{Step 4: Computing the Radon--Nikodym derivative of $T_\psi^* \mu$}

\begin{lemma}\label{lem:RN-shifted}
Let $\rho_\psi = d(T_\psi^* \mu_{\mathrm{GFF}})/d\mu_{\mathrm{GFF}}$ denote the Cameron--Martin density from Lemma \ref{lem:CM}. Then $T_\psi^* \mu$ is absolutely continuous with respect to $\mu_{\mathrm{GFF}}$ with Radon--Nikodym derivative
\[
\frac{d(T_\psi^* \mu)}{d\mu_{\mathrm{GFF}}}(\phi) = F(\phi + \psi) \cdot \frac{1}{\rho_\psi(\phi)} \cdot \rho_\psi(\phi) = F(\phi + \psi) \cdot \frac{\rho_\psi(\phi)}{\int F(\phi'+\psi)\rho_\psi(\phi') d\mu_{\mathrm{GFF}}(\phi')}.
\]
More precisely, $T_\psi^* \mu$ has density
\begin{equation}\label{eq:RN-shifted-final}
\frac{d(T_\psi^* \mu)}{d\mu_{\mathrm{GFF}}}(\phi) = \frac{F(\phi + \psi)}{\int F(\phi' + \psi) \, d\mu_{\mathrm{GFF}}(\phi')}.
\end{equation}
Wait—we need to re-derive this more carefully.
\end{lemma}

We redo this computation carefully.

\begin{lemma}\label{lem:RN-final}
Under the notation above:
\begin{enumerate}
\item[(a)] $T_\psi^* \mu$ is absolutely continuous with respect to $\mu_{\mathrm{GFF}}$.
\item[(b)] The Radon--Nikodym derivative is strictly positive $\mu_{\mathrm{GFF}}$-almost surely.
\item[(c)] $\mu \ll T_\psi^* \mu$ (i.e., $\mu$ is absolutely continuous with respect to $T_\psi^* \mu$).
\end{enumerate}
Hence $\mu$ and $T_\psi^* \mu$ are equivalent.
\end{lemma}

\begin{proof}
We use the change-of-variables formula for pushforward measures.

For any bounded measurable function $f: \mathcal{D}'(\mathbb{T}^3) \to \mathbb{R}$:
\begin{align*}
\int f \, d(T_\psi^* \mu) &= \int f(T_\psi(\phi)) \, d\mu(\phi) = \int f(\phi + \psi) \, d\mu(\phi) \\
&= \int f(\phi + \psi) F(\phi) \, d\mu_{\mathrm{GFF}}(\phi).
\end{align*}

Now apply the Cameron--Martin change of variables $\phi \mapsto \phi - \psi$ (i.e., replace $\phi$ by $\phi - \psi$ in the integral). Since $\psi \in \mathcal{H}$ (the Cameron--Martin space), this change of variables is valid and:
\[
\int f(\phi + \psi) F(\phi) \, d\mu_{\mathrm{GFF}}(\phi) = \int f(\phi) F(\phi - \psi) \rho_{-\psi}(\phi) \, d\mu_{\mathrm{GFF}}(\phi),
\]
where $\rho_{-\psi}(\phi) = d(T_{-\psi}^* \mu_{\mathrm{GFF}})/d\mu_{\mathrm{GFF}}(\phi)$ is the Cameron--Martin density for the shift by $-\psi$.

More precisely: the measure $\mu_{\mathrm{GFF}}$ satisfies $d\mu_{\mathrm{GFF}}(\phi - \psi) = \rho_{-\psi}(\phi) d\mu_{\mathrm{GFF}}(\phi)$, where
\[
\rho_{-\psi}(\phi) = \exp\!\left(-\langle (-\Delta + m^2)\psi, \phi\rangle + \frac{1}{2}\langle (-\Delta + m^2)\psi, \psi\rangle\right).
\]

Therefore:
\[
\int f \, d(T_\psi^* \mu) = \int f(\phi) \underbrace{F(\phi - \psi) \cdot \rho_{-\psi}(\phi)}_{\tilde{F}(\phi)} \, d\mu_{\mathrm{GFF}}(\phi).
\]

This shows:
\begin{equation}\label{eq:RN-Tpsi-mu}
\frac{d(T_\psi^* \mu)}{d\mu_{\mathrm{GFF}}}(\phi) = F(\phi - \psi) \cdot \rho_{-\psi}(\phi) = \tilde{F}(\phi).
\end{equation}

We need to verify:
\begin{itemize}
\item[(i)] $\tilde{F}(\phi) \geq 0$ $\mu_{\mathrm{GFF}}$-a.s.: This holds since $F(\phi - \psi) \geq 0$ and $\rho_{-\psi}(\phi) = e^{(\cdots)} > 0$.
\item[(ii)] $\tilde{F}(\phi) > 0$ $\mu_{\mathrm{GFF}}$-a.s.: This holds since $F(\phi - \psi) > 0$ $\mu_{\mathrm{GFF}}$-a.s.\ (by Lemma \ref{lem:interaction} with $\psi$ replaced by $-\psi$) and $\rho_{-\psi} > 0$ everywhere.
\item[(iii)] $\int \tilde{F} \, d\mu_{\mathrm{GFF}} = 1$: Indeed, $\int f \, d(T_\psi^* \mu) = \int f \tilde{F} \, d\mu_{\mathrm{GFF}}$ and taking $f = 1$ gives $1 = \int \tilde{F} \, d\mu_{\mathrm{GFF}}$.
\item[(iv)] $\tilde{F} \in L^1(\mu_{\mathrm{GFF}})$: This follows from (iii).
\end{itemize}

So $T_\psi^* \mu$ is absolutely continuous with respect to $\mu_{\mathrm{GFF}}$ with density $\tilde{F} > 0$ a.s.

Now, since $\mu \ll \mu_{\mathrm{GFF}}$ with density $F > 0$ a.s., and $T_\psi^* \mu \ll \mu_{\mathrm{GFF}}$ with density $\tilde{F} > 0$ a.s., the sets of $\mu$-null sets and $T_\psi^* \mu$-null sets are both equal to the sets of $\mu_{\mathrm{GFF}}$-null sets. Therefore $\mu \sim T_\psi^* \mu$.

More precisely: For any Borel set $A$,
\begin{align*}
\mu(A) = 0 &\iff \int_A F \, d\mu_{\mathrm{GFF}} = 0 \iff \mu_{\mathrm{GFF}}(A \cap \{F > 0\}) = 0 \\
&\iff \mu_{\mathrm{GFF}}(A) = 0 \quad (\text{since } F > 0 \text{ a.s.}) \\
&\iff \int_A \tilde{F} \, d\mu_{\mathrm{GFF}} = 0 \quad (\text{since } \tilde{F} > 0 \text{ a.s.}) \\
&\iff (T_\psi^* \mu)(A) = 0.
\end{align*}

This proves $\mu$ and $T_\psi^* \mu$ have the same null sets, i.e., they are equivalent.
\end{proof}

\section{Verification Protocol}

We now verify the proof line by line.

\textbf{Step 1 check}: The representation \eqref{eq:RN} uses the Barashkov--Gubinelli construction \cite{BG2020}, which rigorously constructs the $\Phi^4_3$ measure as a probability measure on $H^{-1/2-\epsilon}$ absolutely continuous with respect to the GFF. This is a theorem (not an assumption). The density $F = e^{-V}/Z > 0$ a.s.\ follows because $V < \infty$ a.s.\ (the renormalized interaction is finite a.s., which is part of the construction).

\textbf{Step 2 check (Lemma \ref{lem:CM})}: The Cameron--Martin theorem is a classical result. We verified: (a) $\mu_{\mathrm{GFF}}$ is Gaussian on a separable Banach space; (b) $\psi \in C^\infty \subset H^1 = \mathcal{H}$; (c) the Cameron--Martin norm is finite. The formula \eqref{eq:CM-RN} is the standard Cameron--Martin formula for the GFF. The positivity of the density follows from $e^{(\cdots)} > 0$.

\textbf{Step 3 check (Lemma \ref{lem:interaction})}: The expansion of $V(\phi+\psi)$ uses the Wick product formula $:(\phi+\psi)^n: = \sum_k \binom{n}{k}\psi^{n-k}:\phi^k:$, which holds because the Wick product is defined with respect to the covariance of $\phi$ (not $\phi+\psi$), so $\psi$ is deterministic and can be factored out. Each term $\int \psi^{4-k}:\phi^k: dx$ is in the $k$-th Wiener chaos. By Nelson's hypercontractivity theorem \cite{Nelson1973} (stated precisely: elements of the $n$-th Wiener chaos are in $L^p$ for all $p < \infty$), each term is in $L^p$ for all finite $p$. The finiteness of $e^{\pm\Delta V}$ in $L^1$ follows from the fact that Gaussian chaoses have Gaussian-like tails.

\textbf{Step 4 check (Lemma \ref{lem:RN-final})}: The change of variables $\phi \to \phi - \psi$ under $\mu_{\mathrm{GFF}}$ is valid by the Cameron--Martin theorem (since $\psi \in \mathcal{H}$). Formula \eqref{eq:RN-Tpsi-mu} is derived by direct computation. The positivity $\tilde{F} > 0$ a.s.\ uses $F > 0$ a.s.\ from Step 1.

\textbf{Logical structure}: No circular reasoning. The chain is: GFF structure $\to$ Cameron--Martin $\to$ interaction expansion $\to$ equivalence.

\textbf{Degenerate cases}: If $\psi = 0$, then $T_\psi = \text{id}$ and $T_\psi^* \mu = \mu$, which is trivially equivalent to $\mu$. The proof handles this: $\Delta V = 0$, $\rho_{-\psi} = 1$, $\tilde{F} = F$.

\textbf{Conclusion}: The answer to the question is \textbf{Yes}: the measures $\mu$ and $T_\psi^* \mu$ are equivalent for every smooth $\psi \in C^\infty(\mathbb{T}^3)$, including those not identically zero.

\begin{thebibliography}{99}

\bibitem{BG2020}
N.\ Barashkov and M.\ Gubinelli,
\textit{A variational method for $\Phi^4_3$},
Duke Math.\ J.\ \textbf{169} (2020), no.\ 17, 3339--3415.

\bibitem{DaPratoZabczyk}
G.\ Da Prato and J.\ Zabczyk,
\textit{Stochastic Equations in Infinite Dimensions},
Second edition, Cambridge University Press, 2014.

\bibitem{Hairer2014}
M.\ Hairer,
\textit{A theory of regularity structures},
Invent.\ Math.\ \textbf{198} (2014), no.\ 2, 269--504.

\bibitem{MW2017}
J.-C.\ Mourrat and H.\ Weber,
\textit{The dynamic $\Phi^4_3$ model comes down from infinity},
Comm.\ Math.\ Phys.\ \textbf{356} (2017), no.\ 3, 673--753.

\bibitem{Nelson1973}
E.\ Nelson,
\textit{The free Markoff field},
J.\ Funct.\ Anal.\ \textbf{12} (1973), 211--227.

\bibitem{Simon}
B.\ Simon,
\textit{The $P(\phi)_2$ Euclidean (Quantum) Field Theory},
Princeton University Press, 1974.

\end{thebibliography}

\end{document}
